% Third numerical example: fluidic pinball and comparison with alternative ROMs.
% Recommended packages in the main manuscript:
% \usepackage{amsmath,amssymb,booktabs,graphicx,threeparttable}
% The image paths below are relative to the Pinball project root.
% In a single-column manuscript, table* and figure* may be replaced by
% table and figure, respectively.

```latex
\subsection{Fluidic pinball}
\label{subsec:fluidic_pinball}

The fluidic pinball is considered as a third numerical test of RAL-MoE-ROM
over a parameter range spanning steady, Hopf-onset, and periodic wake
dynamics. In addition to evaluating the regime-local specialists and their
overlap coupling, this example compares the proposed reduced dynamics with a
classical POD--Galerkin/Pressure--Poisson ROM\cite{holmes1996turbulence} and with a global
POD-DeepONet model\cite{lu2021learning}. 

The high-fidelity model solves the two-dimensional incompressible
Navier--Stokes equations around three fixed circular cylinders. The cylinder
diameter and free-stream velocity are $D=U_\infty=1$, giving
\[
Re=\frac{U_\infty D}{\nu}=\frac{1}{\nu}.
\]
The cylinders have radius $0.5$ and form an equilateral triangle of side
length $1.5$, with centres
$(-1.2990381,0)$, $(0,0.75)$, and $(0,-0.75)$ in the computational domain
$[-6,20]\times[-6,6]$. A uniform velocity $(1,0)$ is prescribed at the
inlet, slip conditions are imposed on the upper and lower boundaries, and
no-slip conditions are imposed on the cylinder surfaces. The outlet
pressure is fixed.

The full-order simulations are performed with OpenFOAM~13 on a
finite-volume mesh containing 70,194 cells. The CFD time step is
$\delta t_{\mathrm{CFD}}=0.005$, and snapshots are stored every $0.25$
physical time units. The database contains 131 complete Reynolds-number trajectories over
$10\le Re\le60$, partitioned into 94 training, 19 validation, and 18
held-out test trajectories. The same Reynolds-number trajectory retains the
same split role when it belongs to two overlapping regime-local datasets.
Table~\ref{tab:pinball_data_contracts} summarizes the three specialist
domains, their trajectory-level partitions, and the corresponding POD
dimensions. 

\begin{table*}[htbp]
\centering
\caption{Data partitions and local POD dimensions for the fluidic-pinball
case.}
\label{tab:pinball_data_contracts}

\small
\renewcommand{\arraystretch}{1.10}
\setlength{\tabcolsep}{8pt}

\begin{tabular}{@{}lcccc@{}}
\toprule
Regime
& Reynolds-number range
& Train/validation/test
& $r_u$
& $r_p$ \\
\midrule

Steady
& $10.0\le Re\le19.0$
& $37/8/7$
& 4
& 3 \\

Hopf
& $16.6\le Re\le25.0$
& $39/8/7$
& 5
& 5 \\

Periodic
& $21.25\le Re\le60.0$
& $47/9/9$
& 17
& 16 \\

\bottomrule
\end{tabular}
\end{table*}

The corresponding overlap regions are
$16.6\le Re\le19.0$ for $S$--$H$ and
$21.25\le Re\le25.0$ for $H$--$P$. The proposed specialists follow the
common regime-local formulation introduced in Section~3, and the outer
routing and T2-C combination follow Section~4. The field-error measures are
those introduced in the common numerical evaluation protocol.

We first compare the proposed regime-local dynamics with the
POD--Galerkin ROM under matched local representations and
prediction horizons. For these comparisons, both models use the same
train-only local POD spaces and are evaluated from identical initial
histories. The Periodic comparison uses the cadence aligned with the
specialist history, corresponding to $\Delta t=1$.

\begin{table*}[htbp]
\centering
\caption{Matched long-horizon comparison between the classical
POD--Galerkin ROM and the proposed regime-local
specialists. Errors are percentages.}
\label{tab:pinball_galerkin_matched}

\small
\renewcommand{\arraystretch}{1.10}
\setlength{\tabcolsep}{5pt}

\begin{tabular}{@{}llcccc@{}}
\toprule
Regime and horizon
& Model
& $E_u$
& $E_p$
& $E_{\mathrm{joint}}$
& Divergent windows \\
\midrule

Steady, $K=56$
& POD-Galerkin
& --
& --
& --
& $448/448$ \\

& Proposed specialist
& 0.2247
& 0.7593
& 0.9840
& 0 \\

\midrule

Hopf matched subset, $K=56$
& POD-Galerkin
& 0.4378
& 48.50
& 48.94
& $141/320$ \\

& Proposed specialist
& 0.0995
& 1.986
& 2.085
& 0 \\

\midrule

Periodic, $K=32$
& POD-Galerkin
& 1.641
& 41.56
& 43.20
& 0 \\

& Proposed specialist
& 0.2897
& 0.5518
& 0.8415
& 0 \\

\midrule

Periodic core, $K=56$
& POD-Galerkin
& 2.555
& 39.36
& 41.91
& 0 \\

& Proposed specialist
& 0.6222
& 1.050
& 1.672
& 0 \\

\bottomrule
\end{tabular}
\end{table*}

Table~\ref{tab:pinball_galerkin_matched} shows clear differences between
the classical POD-Galerkin reduced dynamics and the
proposed regime-local specialists under autonomous rollout. In the Steady
regime, the POD-Galerkin ROM trajectories become non-finite or exceed the
prescribed divergence criterion before a complete $K=56$ score can be formed,
whereas the proposed specialist remains finite and achieves a joint error
below $1\%$. The Hopf comparison reveals a similar limitation of the classical reduced
dynamics. Although finite field values are obtained for the matched windows,
a substantial fraction crosses the divergence threshold, and the resulting
error is dominated by the pressure component. For the available matched
finite windows, the POD-Galerkin model gives a joint error of $48.94\%$,
whereas the proposed specialist achieves $2.085\%$ without divergence. The Periodic regime enables a direct quantitative comparison because both
models remain finite over the prescribed rollout horizons. At $K=32$, the
proposed specialist reduces the joint error from $43.20\%$ to $0.8415\%$;
over the longer $K=56$ Periodic-core rollout, the error decreases from
$41.91\%$ to $1.672\%$. The dominant difference originates from the pressure
prediction, for which th ePOD-Galerkin model exhibits errors of approximately
$40\%$, whereas the proposed specialist remains close to $1\%$. These results
indicate that the projected reduced dynamics with the standard
Pressure--Poisson closure are insufficient for accurate long-horizon
autonomous prediction in this problem, even when the underlying POD
representation is fixed.



A complementary comparison is provided by the global POD-DeepONet. On
4,228 held-out one-step pairs, the global model gives
$E_u=1.5921\%$ and $E_p=23.3447\%$. The finite one-step predictions do not
persist under recursive application: only one of the 28 unique held-out
trajectories remains finite over autonomous rollout, whereas the remaining
trajectories develop non-finite states. The one-step and autonomous results
therefore describe fundamentally different properties of the learned
operator. In particular, finite one-step approximation error does not imply
that repeated application defines a stable reduced dynamical evolution.



We next examine the output-level combination of adjacent regime-local
specialists. Table~\ref{tab:pinball_boundary_comparison} reports the two
individual specialist predictions, direct blending based on the E2
probabilities, T2-C, and the corresponding window-wise convex oracle. The
$S$--$H$ comparison uses the Steady and Hopf specialists, whereas the
$H$--$P$ comparison uses the Periodic and Hopf specialists.

\begin{table*}[htbp]
\centering
\caption{Prediction errors for the adjacent-specialist combinations in the
fluidic-pinball overlap regions. All entries are
$E_{\mathrm{joint}}$ (\%).}
\label{tab:pinball_boundary_comparison}

\small
\renewcommand{\arraystretch}{1.10}
\setlength{\tabcolsep}{6pt}

\begin{tabular}{@{}llccccc@{}}
\toprule
Overlap
& Horizon
& Primary specialist
& Hopf specialist
& E2-probability blend
& T2-C
& Convex oracle \\
\midrule

$S$--$H$
& $K=24$
& 0.6014
& 3.1925
& 1.7511
& 0.5897
& 0.5889 \\

$S$--$H$
& $K=56$
& 0.6927
& 3.1193
& 1.7714
& 0.6767
& 0.6763 \\

\midrule

$H$--$P$
& $K=24$
& 0.0990
& 0.5146
& 0.2972
& 0.0945
& 0.0928 \\

$H$--$P$
& $K=56$
& 0.1028
& 0.5346
& 0.3085
& 0.0978
& 0.0955 \\

\bottomrule
\end{tabular}
\end{table*}




For the two overlap regions, T2-C consistently improves upon the dominant
individual specialist or hard fusion strategy. In the $S$--$H$ overlap, the
Steady specialist provides the more accurate individual prediction, while
direct use of the E2 probabilities as fusion weights degrades the result.
T2-C reduces the $K=56$ joint error from $0.6927\%$ to $0.6767\%$, approaching
the convex-oracle value of $0.6763\%$. The improvement is modest in magnitude, but its
consistency indicates that the adjacent Hopf prediction contains a useful
corrective component that can be exploited without replacing the dominant
Steady contribution. A similar
trend is observed in the $H$--$P$ overlap, where T2-C reduces the joint error
from $0.1028\%$ to $0.0978\%$ at $K=56$. 



The field reconstructions provide spatial context for the long-horizon
errors. Figure~\ref{fig:pinball_hopf_velocity_errors} shows a representative
Hopf validation trajectory, whereas
Fig.~\ref{fig:pinball_periodic_velocity_errors} shows a held-out Periodic
trajectory. In each case, the physical velocity field and its pointwise
absolute error are shown at an intermediate and terminal prediction time.

\begin{figure*}[htbp]
\centering

\includegraphics[width=0.90\textwidth]
{DDPDCNN/4/RAL_MOE_ROM_METHODS_OVERLEAF_READY_20260726/drafts/cylinder_wake_20260726/figures/SH_3.png}

\vspace{0.5em}

\includegraphics[width=0.90\textwidth]
{DDPDCNN/4/RAL_MOE_ROM_METHODS_OVERLEAF_READY_20260726/drafts/cylinder_wake_20260726/figures/SH_3_2.png}

\caption{Representative Hopf-specialist prediction at $Re=19.9$ on a
validation trajectory. The upper panel compares the CFD reference and the
proposed prediction at $K=28$ and $K=56$, while the lower panel shows the
corresponding pointwise absolute velocity errors.}
\label{fig:pinball_hopf_velocity_errors}
\end{figure*}

\begin{figure*}[htbp]
\centering

\includegraphics[width=0.90\textwidth]
{DDPDCNN/4/RAL_MOE_ROM_METHODS_OVERLEAF_READY_20260726/drafts/cylinder_wake_20260726/figures/PH_3.png}

\vspace{0.5em}

\includegraphics[width=0.90\textwidth]
{DDPDCNN/4/RAL_MOE_ROM_METHODS_OVERLEAF_READY_20260726/drafts/cylinder_wake_20260726/figures/PH_3_2.png}

\caption{Representative held-out Periodic-specialist prediction at $Re=43$.
The upper panel compares the CFD reference and the proposed prediction at
$K=28$ and $K=56$, while the lower panel shows the corresponding pointwise
absolute velocity errors.}
\label{fig:pinball_periodic_velocity_errors}
\end{figure*}

In the Hopf example, the velocity discrepancy remains localized primarily
around the separated shear layers and the near-wake region. The error does
not exhibit substantial downstream amplification over the considered rollout
horizon. In the Periodic example, the spatial error extends along the vortex
street as the prediction horizon increases, while the alternating wake
structure remains preserved. This error pattern is consistent with the
accumulation of phase displacement in coherent structures rather than a
transition toward an incorrect dynamical regime. 

Overall, the fluidic-pinball example highlights three aspects of the proposed
framework. The matched comparison with the POD--Galerkin ROM indicates that
the classical projected reduced dynamics with Pressure--Poisson closure do
not provide sufficient long-horizon autonomous accuracy for this problem,
particularly in the pressure prediction, whereas the proposed regime-local
specialists maintain substantially lower errors. The global POD-DeepONet
comparison further shows that accurate one-step operator approximation does
not necessarily translate into stable recursive reduced dynamics. Finally,
the overlap results indicate that adjacent regime-local specialists contain
complementary predictive information that can be exploited through
output-level adaptive fusion. The available $H$--$P$ evaluation provides
additional evidence of such complementarity under the adopted evaluation
setting.

